\documentclass[12pt,a4paper]{article}

% Idioma y codificación
\usepackage[spanish,es-tabla]{babel}
\usepackage[T1]{fontenc}
\usepackage[utf8]{inputenc}

% Formato y utilidades
\usepackage[a4paper,margin=2.5cm]{geometry}
\usepackage{graphicx}
\graphicspath{{anexos/}{figuras/}{imagenes/}{img/}}
\usepackage{float}      % [H]
\usepackage{hyperref}
\usepackage{booktabs}
\hypersetup{
  colorlinks=true,
  linkcolor=black,
  urlcolor=blue,
  citecolor=black
}

% Bibliografía
\usepackage[style=ieee,backend=biber]{biblatex}
\addbibresource{referencias.bib}

\begin{document}

% Portada simple
\begin{titlepage}
  \centering
  {\Large Universidad Tecnol\'ogica del Uruguay (UTEC)\par}
  \vspace{0.5cm}
  {\large Proyecto Integrador — Brazo rob\'otico did\'actico\par}
  \vspace{2cm}
  {\LARGE Evidencias del proyecto\par}
  \vspace{1.5cm}
  \begin{tabular}{rl}
    Integrantes: & Hector Pereira, Mateo Lecuna, Priscila Rossi \\
    Docente/s:   & Marcelo Diaz, Christian Campero \\
    Fecha:       & \today
  \end{tabular}
  \vfill
\end{titlepage}

\tableofcontents
\newpage

\section{Resumen}
% Breve descripción del objetivo y alcance. Qué contiene este documento.

\section{Metodología y cronograma}

\section{Evidencias de diseño y fabricación}


\begin{table}[H]
\centering
\caption{Evaluación comparativa de mecanismos de actuación}
\begin{tabular}{lccccc}
\toprule
\textbf{Atributo} &\textbf{Peso} &\textbf{Servom.} &\textbf{Eslab.} &\textbf{Actuad. lin.} &\textbf{Pistones hidr.} \\
\midrule
Distribución de peso & 1 & 6  & 8  & 9 & 10 \\
Simplicidad          & 3 & 10 & 6  & 5 & 4  \\
Costo                & 3 & 9  & 8  & 7 & 7  \\
Reparabilidad        & 2 & 9  & 7  & 6 & 5  \\
Durabilidad          & 2 & 8  & 5  & 4 & 3  \\
Cap. de carga      & 2 & 7  & 7  & 8 & 9  \\
Velocidad            & 1 & 8  & 8  & 5 & 4  \\
Precisión            & 1 & 8  & 7  & 6 & 5  \\
\midrule
\textbf{Total ponderado}       &    & \textbf{127} & \textbf{103} & \textbf{92} & \textbf{86} \\
\bottomrule
\end{tabular}
\end{table}

\begin{table}[H]
\centering
\caption{Matriz de selección entre servos MG996R y SG90}
\begin{tabular}{lccc}
\toprule
\textbf{Atributo} & \textbf{Peso} & \textbf{Servo MG996R} & \textbf{Servo SG90} \\
\midrule
Fuerza (torque)        & 4 & 10 & 3 \\
Tamaño (dimensiones)   & 2 & 6  & 9 \\
Peso (masa)            & 2 & 5  & 9 \\
Precio                 & 2 & 7  & 9 \\
Durabilidad / materiales & 3 & 9 & 5 \\
Precisión              & 2 & 7  & 6 \\
Velocidad (s/60°)      & 1 & 6  & 9 \\
Corriente (menor=mejor) & 2 & 4  & 8 \\
\midrule
\textbf{Total ponderado} &   & \textbf{131} & \textbf{118} \\
\bottomrule
\end{tabular}
\end{table}


\begin{table}[H]
\centering
\caption{Matriz de selección: MDF vs Impresión 3D vs Híbrido (MDF+PLA)}
\begin{tabular}{lcccc}
\toprule
\textbf{Atributo} & \textbf{Peso} & \textbf{MDF} & \textbf{PLA} & \textbf{MDF+PLA} \\
\midrule
Tiempo de fabricación (menor=mejor) & 3 & 9 & 3 & 8 \\
Costo de materiales (menor=mejor)   & 2 & 9 & 5 & 8 \\
Facilidad de iteración/retrabajo    & 3 & 10 & 4 & 9 \\
Precisión / tolerancias              & 2 & 6 & 8 & 8 \\
Comportamiento a la fricción         & 3 & 3 & 8 & 8 \\
Rigidez / resistencia estructural    & 3 & 7 & 6 & 8 \\
Peso total (menor=mejor)             & 1 & 6 & 7 & 7 \\
Complejidad de ensamblaje (menor=mejor) & 1 & 8 & 6 & 6 \\
Disponibilidad de recursos (UTEC)    & 1 & 10 & 6 & 9 \\
Acabado / ajuste final               & 1 & 6 & 7 & 7 \\
\midrule
\textbf{Total ponderado}             &    & \textbf{147} & \textbf{115} & \textbf{160} \\
\bottomrule
\end{tabular}
\end{table}

\begin{table}[H]
\centering
\caption{Matriz de selección de fuente de poder}
\begin{tabular}{lccccc}
\toprule
\textbf{Atributo} & \textbf{Peso} & \textbf{PC} & \textbf{LED} & \textbf{Genéricas} & \textbf{Fuente de lab.} \\
\midrule
Corriente continua y picos         & 3 & 9 & 7 & 6 & 8 \\
Disponibilidad + costo             & 2 & 10 & 7 & 8 & 3 \\
Facilidad de integración/ensamble  & 2 & 6 & 9 & 4 & 5 \\
Regulación y rizado                & 2 & 8 & 6 & 7 & 9 \\
Protecciones (OCP/OVP/OTP, etc.)   & 2 & 9 & 6 & 7 & 9 \\
Portabilidad / formato             & 1 & 5 & 7 & 6 & 3 \\
Ruido eléctrico/EMI                & 1 & 7 & 6 & 6 & 8 \\
Escalabilidad (margen futuro)      & 2 & 9 & 6 & 5 & 7 \\
Confiabilidad global               & 2 & 8 & 7 & 5 & 8 \\
\midrule
\textbf{Total ponderado}           &   & \textbf{139} & \textbf{116} & \textbf{102} & \textbf{117} \\
\bottomrule
\end{tabular}
\end{table}

\begin{table}[H]
\centering
\caption{Matriz de selección de microcontrolador}
\begin{tabular}{lcccc}
\toprule
\textbf{Atributo} & \textbf{Peso} & \textbf{UNO} & \textbf{NANO} & \textbf{ESP32} \\
\midrule
Disponibilidad en UTEC                     & 3 & 10 & 9  & 5  \\
Integración con servos/5 V                 & 3 & 9  & 9  & 6  \\
Ecosistema y soporte en curso              & 2 & 10 & 9  & 7  \\
Costo                                      & 2 & 8  & 9  & 8  \\
Tamaño                                     & 1 & 7  & 10 & 8  \\
Rendimiento (CPU/RAM)                      & 1 & 6  & 6  & 10 \\
Conectividad (UART extra/WiFi/BLE)         & 1 & 3  & 3  & 10 \\
Consumo (menor = mejor)                    & 1 & 8  & 9  & 5  \\
Robustez en banco (taller/educativo)       & 1 & 8  & 7  & 7  \\
Expansión de E/S y temporizadores          & 1 & 8  & 7  & 9  \\
\midrule
\textbf{Total ponderado}                 &   & \textbf{133} & \textbf{132} & \textbf{112} \\
\bottomrule
\label{tab:matriz_micro}
\end{tabular}
\end{table}



\subsection{Materiales y métodos}




\subsection{Galgas y tolerancias}
% Descripción de pruebas y cómo se midió. Incluir figuras si corresponde.

\subsection{Articulaciones y prototipos}
% Versiones, cambios y fotos de las piezas.

\section{Electrónica}
% Esquema, distribución de potencia, cableado. Fotos/figuras opcionales.

\section{Programación}
% Modos de operación, pruebas básicas, logs o capturas.

\section{Pruebas y resultados}
% Qué se probó, evidencias fotográficas o capturas.

\section{Documentación generada}
% Manuales, planos, listas, enlaces a repositorio.

\section{Conclusiones}
% Hallazgos clave y posibles mejoras.

\appendix
\section{Inspiraciones y referencias de mercado}
% Enlaces o capturas. Referido en el cuerpo principal.


\end{document}